\documentclass[twocolumn,11pt]{article}
\usepackage[margin=0.85in]{geometry}
\usepackage{booktabs}
\usepackage{amsmath,amssymb}
\usepackage{hyperref}
\usepackage{natbib}
\usepackage{enumitem}
\hypersetup{colorlinks=true,linkcolor=blue,citecolor=blue,urlcolor=blue}

\newcommand{\lcdm}{$\Lambda$CDM}
\newcommand{\Msun}{M_\odot}
\newcommand{\einf}{\varepsilon_\infty}
\newcommand{\Ndet}{N_{\rm det}}
\newcommand{\Csub}{C_{\rm sub}}

\begin{document}

\title{Erratum to ISF Papers VI, VII and VIII:\\
a missing Jacobian in the per-lens perturber count\\[0.3em]
\large Every detectable-perturber count in the three papers is high by
a factor 1.4, every Poisson significance by $\sqrt{1.4}$; the joint
weak-lensing fit, the profiles and the falsification logic stand}

\author{Christopher K.\ Merrill\\[0.3em]
{\small With computational collaboration from Claude Fable 5.1 (Anthropic)}}
\date{October 11, 2026}

\maketitle

\begin{abstract}
Paper VI of the Independent Structure Formation (ISF) series counts the
other-branch galaxies detectable per strong-lens sightline by
integrating a stellar mass function, written per dex of stellar mass,
over the logarithm of halo mass. The conversion between the two carries
the Jacobian $d\log M_\ast/d\log M_{\rm halo}$, which for the paper's
single power-law abundance matching is the constant $1/1.4$, and the
script omitted it. Every quantity downstream of that integral is
therefore high by exactly 1.4: Paper VI's $1.08$ detectable perturbers
per lens at $N=6$ is $0.77$; Paper VII's calibration constant
$\Csub=0.514$ is $0.367$ and its per-lens counts at $N=2$--10 fall from
$0.51$--$1.20$ to $0.37$--$0.86$; Paper VII's fiducial Euclid
expectation of 4,050 detections is 2,893 and its all-ISF falsification
threshold of 1,844 is 1,307; Paper VIII's wide-survey expectation of
4,463--5,955 is 3,188--4,253, its threshold of 2,037--2,733 is
1,444--1,939, its DR1 null exclusion of $\geq24\sigma$ is
$\geq20\sigma$, and its 312-system in-principle test of 12.6--16.8
expected perturbers at 5.0--5.8$\sigma$ is 9.0--12.0 at 4.2--4.9$\sigma$.
Nothing is strengthened and nothing is fatal; the series' falsification
statements are restated here with the corrected numbers. One
qualitative statement does change: in the ISF-true forecast scenario the
weaker Euclid term no longer collapses the posterior onto the $N=6$
that generated the mock, which splits between $N=6$ and $N=8$. The
audit also found that Paper VI's text quotes Schechter parameters its
script does not use; the published numbers come from the script's
parameters, and the text is corrected to them. Paper VII's joint
weak-lensing plus substructure fit ($N=10$, $k_t^{\rm phys}=0.52\,h\,{\rm
Mpc}^{-1}$), Paper VI's convergence profiles and bound-mass table, Paper
IX, and Paper X's own counts are untouched. Every number here is
printed by a committed script.
\end{abstract}

\section{The error}
\label{sec:error}

Paper VI \citep{paper6} counts the detectable other-branch galaxies
per strong-lens sightline as
\begin{equation}
\begin{split}
\Ndet=\int_{10}^{12}(N-1)\bigl[1-\varepsilon(k_M)\bigr]\,
\Phi\bigl(M_\ast(M_{\rm halo})\bigr)\\
\times\,V\;d\log_{10}M_{\rm halo},
\end{split}
\label{eq:integral}
\end{equation}
with $\Phi$ the stellar mass function per dex of \emph{stellar} mass,
evaluated at the abundance-matched stellar mass
$M_\ast=10^{10.7}(M_{\rm halo}/10^{12}\,\Msun)^{1/1.4}$, and the sum
taken over $d\log M_{\rm halo}$. Paper VI's equation for $dn_{\rm ISF}/d\log_{10}M$ calls
$\Phi_{\rm gal}(M)$ ``the galaxy halo mass function (converted from
the stellar mass function via the stellar-to-halo mass relation)''.
That conversion is
\begin{equation}
\frac{dn}{d\log M_{\rm halo}}=\frac{dn}{d\log M_\ast}\,
\frac{d\log M_\ast}{d\log M_{\rm halo}}
=\Phi(M_\ast)\,\frac{1}{1.4},
\end{equation}
and the script that produced the published numbers omits the second
factor. Because the matching relation is a single power law the factor
is a constant, so every quantity downstream of Eq.~\ref{eq:integral}
is high by exactly 1.4 and every Poisson significance built on it by
$\sqrt{1.4}=1.18$. Reproducing Paper VI's functions with its exact
constants gives the published $1.0806$ per lens at $N=6$,
$\einf=0.58$ (and $1.80$, $2.57$ at $\einf=0.30$, $0$) to three parts
in $10^4$; with the Jacobian the same code gives $0.772$, $1.287$,
$1.838$, a ratio of $1.3996$.

The omission was found during the peer review of Paper X
\citep{paper10} on 2026-10-10, confirmed the same day against Paper
VI's text and script, and declared in Paper X as its concession~6.
Paper X gates on Paper VI as published and its own perturber counts use
a mapping with unit Jacobian, so they do not change (Section~\ref{sec:paperx}).

\section{The Schechter parameters}
\label{sec:schechter}

The audit found a second discrepancy that the erratum must settle.
Paper VI's text quotes a single Schechter function with
$\Phi^\ast=0.012\;{\rm Mpc^{-3}\,dex^{-1}}$,
$\log_{10}M_\ast^{\rm char}=10.7$ and $\alpha_S=-1.05$, attributed to
\citet{baldry12}. The script that produced the published numbers uses
$\Phi^\ast=3.96\times10^{-3}$, $\log_{10}M_\ast^{\rm char}=10.7$,
$\alpha_S=-1.16$. Neither is \citet{baldry12} as published, which is a
double Schechter with $\log_{10}M_\ast^{\rm char}=10.66$,
$\Phi_1^\ast=3.96\times10^{-3}$, $\alpha_1=-0.35$,
$\Phi_2^\ast=0.79\times10^{-3}$, $\alpha_2=-1.47$. Table~\ref{tab:schechter}
gives the per-lens count each set produces. The text's parameters
would have given 2.65 per lens, 2.4 times the published value, so the
discrepancy is not cosmetic. The correction is to the text: Paper VI's
Schechter equation should read $\Phi^\ast=3.96\times10^{-3}\;{\rm Mpc^{-3}\,dex^{-1}}$,
$\log_{10}M_\ast^{\rm char}=10.7$, $\alpha_S=-1.16$, a single-Schechter
approximation whose normalisation and characteristic mass follow
\citet{baldry12} and whose low-mass slope is an intermediate value.
The double Schechter as published would give $0.49$ per lens with the
Jacobian, 37 percent below the corrected $0.77$. That is a model change,
not an erratum; it is reported here as a check and not adopted, and a
future version of the series may take it up.

\begin{table*}
\centering
\caption{Detectable perturbers per lens at $N=6$, $\einf=0.58$, under
the three Schechter parameter sets, without and with the Jacobian. The
published 1.08 is the script set without it; the corrected value is
the script set with it.}
\label{tab:schechter}
\small
\begin{tabular}{lcc}
\toprule
parameter set & no Jacobian & with Jacobian \\
\midrule
text ($\Phi^\ast$ 0.012, $\alpha_S$ $-1.05$)        & 2.647 & 1.891 \\
script ($3.96\times10^{-3}$, $-1.16$)               & \textbf{1.081} & \textbf{0.772} \\
Baldry et al.\ (2012) double Schechter              & 0.681 & 0.486 \\
\bottomrule
\end{tabular}
\end{table*}

\section{Corrected numbers}
\label{sec:numbers}

Tables~\ref{tab:p6}--\ref{tab:p8} list every number the three papers
quote from the affected chain, published beside corrected. The Paper VII
and VIII values come from rerunning the papers' own forecast scripts on
the corrected per-lens counts, after each rerun first reproduced the
published data product from the published input to a relative
$10^{-9}$. The corrected products are published beside the originals as
\texttt{euclid\_forecast\_v2.npz} and
\texttt{recalibrated\_forecast\_v2.npz}; the originals are not altered.

\begin{table*}
\centering
\caption{Paper VI. $N_{\rm sub}$, the bound-subhalo channel, is
$3.2\times10^{-4}$ per lens and stays negligible. The published
Euclid row used 1.081 rounded; the corrected row is from the script.}
\label{tab:p6}
\small
\setlength{\tabcolsep}{4pt}
\begin{tabular}{lcc}
\toprule
quantity & published & corrected \\
\midrule
$\Ndet$ per lens, $\einf=0.58$ (abundance table)   & 1.08 & 0.77 \\
$\Ndet$ per lens, $\einf=0.30$             & 1.80 & 1.29 \\
$\Ndet$ per lens, $\einf=0.00$             & 2.57 & 1.84 \\
ratio to the CDM baseline 0.43             & $2.5\times$ & $1.8\times$ \\
combined $N_{\rm total}$ per lens (LOS table) & 1.13 & 0.82 \\
combined $f_{\rm LOS}$ / $f_{\rm sub}$     & 98.5 / 1.5\% & 98.0 / 2.0\% \\
LOS-fraction shift from \lcdm              & 31.5 points & 31.0 points \\
Euclid, $\einf=0.58$: per lens / 75,000 / 5\% & 1.081 / 81,075 / 4,054 & 0.772 / 57,904 / 2,895 \\
Euclid, $\einf=0.30$                       & 1.801 / 135,075 / 6,754 & 1.287 / 96,506 / 4,825 \\
abstract: ``$\sim$4,000 at $>50\sigma$''   & 4,054 at 90$\sigma$ & 2,895 at 76$\sigma$ \\
2,500 imaging-suitable lenses              & 2,701 & 1,930 \\
criterion (1): $>$1,000 lenses, $\eta=0.05$ & 10.4$\sigma$ & 8.8$\sigma$ \\
\bottomrule
\end{tabular}
\end{table*}

\begin{table*}
\centering
\caption{Paper VII, Section 6 (75,000 grade-A/B lenses, $\eta=0.05$,
$10^{8.8}\,\Msun$). Unchanged: $\einf(N)$, $f_{\rm GGSL}$, the joint
fit, and the current $P(N)$.}
\label{tab:p7}
\small
\setlength{\tabcolsep}{4pt}
\begin{tabular}{lcc}
\toprule
quantity & published & corrected \\
\midrule
$\Csub$                                       & 0.514 & 0.367 \\
$\Ndet$ per lens, $N=2/3/4$                   & 0.51 / 0.77 / 0.93 & 0.37 / 0.55 / 0.66 \\
$\Ndet$ per lens, $N=6/8/10$                  & 1.08 / 1.15 / 1.20 & 0.77 / 0.82 / 0.86 \\
expected detections, $N=2$ to 10              & 1,929--4,513 & 1,378--3,223 \\
fiducial $N=6$: $n_{\rm exp}$, Poisson $\sigma$ & 4,050, 63.6 & 2,893, 53.8 \\
fractional precision on $\Ndet$               & 1.6\% & 1.9\% \\
ISF-true: $N=2$ / 3 / 4 from $N=6$            & 33.3 / 18.2 / 9.1$\sigma$ & 28.2 / 15.4 / 7.7$\sigma$ \\
ISF-true: $N=8$ / 10 from $N=6$               & 4.2 / 7.3$\sigma$ & 3.6 / 6.1$\sigma$ \\
ISF-true: $P(N=6)$, $P(N=8)$                  & 0.92, 0.08 & 0.47, 0.52 \\
ISF-true joint best fit                       & $N=6$, $k_t=0.79$ & $N=8$, $k_t=1.08$ \\
null: $N=2$ / 6 / 10 excluded at              & 62 / 90 / 95$\sigma$ & 52 / 76 / 80$\sigma$ \\
null: min $\Delta\chi^2$ vs \lcdm             & 3,859 & 2,757 \\
95\% threshold, all ISF ($N\geq2$)            & $<$1,844 & $<$1,307 \\
95\% threshold, $N\geq4$                      & $<$3,358 & $<$2,384 \\
\bottomrule
\end{tabular}
\end{table*}

\begin{table*}
\centering
\caption{Paper VIII (measured Q1 demographics; brackets over the flat,
linear and area efficiency models). Unchanged: the effective
efficiencies 0.050 / 0.041 / 0.038, the stages Q1 497, DR1 14,965,
Wide 110,269, and the 312-system effective sightlines 15.6 / 12.7 /
11.7.}
\label{tab:p8}
\small
\setlength{\tabcolsep}{4pt}
\begin{tabular}{lcc}
\toprule
quantity & published & corrected \\
\midrule
Wide $n_{\rm exp}(N=6)$                      & 4,463--5,955 & 3,188--4,253 \\
Wide 95\% threshold, all ISF                 & 2,037--2,733 & 1,444--1,939 \\
Wide 95\% threshold, $N\geq4$                & 3,706--4,966 & 2,632--3,529 \\
Wide ISF-true $P(N=6)$                       & 0.967--0.999 & 0.635--0.949 \\
Wide null $\sigma(N=2)$                      & 65.2--75.3 & 55.1--63.6 \\
DR1 $n_{\rm exp}(N=6)$                       & 606--808 & 433--577 \\
DR1 95\% threshold, all ISF                  & 257--348 & 180--244 \\
DR1 null $\sigma(N=2)$ / $\sigma(N=6)$       & $\geq24.0$ / $\geq34.8$ & $\geq20.3$ / $\geq29.4$ \\
DR1 ISF-true $P(N=6)$                        & $\sim10^{-3}$ & $\sim10^{-3}$ \\
DR1 $\sigma(4\,{\rm v}\,6)$                  & 3.5--4.1 & 3.0--3.4 \\
Q1 $n_{\rm exp}(N=6)$                        & 20--27 & 14--19 \\
312 systems: $n_{\rm exp}(N=6)$              & 12.6--16.8 & 9.0--12.0 \\
312 systems, zero found: $N\geq6$ excluded at & 5.0--5.8$\sigma$ & 4.2--4.9$\sigma$ \\
312 systems: $N=10$ excluded at              & 5.3--6.1$\sigma$ & 4.5--5.2$\sigma$ \\
312 systems: all $N\geq2$ excluded at        & 3.5--4.0$\sigma$ & 2.9--3.4$\sigma$ \\
\bottomrule
\end{tabular}
\end{table*}

\section{What changes qualitatively}
\label{sec:qualitative}

Two statements do not survive as written; both concern the ISF-true
forecast scenario, in which a mock Euclid count generated at $N=6$ is
added to the current joint fit.

\textbf{Paper VII, Section 6.1.} ``$P(N=6)=0.92$ (up from
$9\times10^{-6}$ currently)'' and ``the current WL-dominated posterior
($N=10$ preferred) shifts decisively to $N=6$'' are wrong after the
correction. With the smaller count the Euclid term separates $N=6$ from
$N=8$ by 3.6$\sigma$ rather than 4.2, which no longer overrides the
current fit's preference for $N=8$--10: the corrected posterior is
$P(N=6)=0.47$, $P(N=8)=0.52$, $P(N=10)=3\times10^{-4}$, with the joint
best fit at $N=8$. The forecast still excludes $N\leq4$ and $N=10$ at
the mock; the claim weakens from ``Euclid collapses the $N$ posterior''
to ``Euclid narrows it to $N=6$--8''.

\textbf{Paper VIII, abstract and Section 4.} ``Only the full wide
survey collapses the posterior onto the truth ($P(N=6)=0.97$--$1.00$)''
becomes $P(N=6)=0.64$--$0.95$: the wide survey still prefers $N=6$,
decisively under the flat efficiency model and no longer decisively
under the linear and area models. The companion statement, that DR1
falsifies but cannot confirm, survives unchanged in form: the DR1 null
exclusion stays above 20$\sigma$ and the DR1 ISF-true $P(N=6)$ stays
near $10^{-3}$.

Everything else is a rescaling. ``DR1 is decisive for falsification''
holds at $\geq20\sigma$; ``the 312 already-modelled Q1 systems suffice
for a first real-data test'' holds at $\geq4.2\sigma$ for $N\geq6$,
which no longer reaches the 5$\sigma$ the abstract quoted, and at
2.9--3.4$\sigma$ for all $N\geq2$.

\section{What is unchanged}
\label{sec:unchanged}

Verified by reading the code and, where a product exists, by byte
comparison after the reruns:
\begin{itemize}[leftmargin=1.4em,itemsep=0.1em]
\item Paper VII's joint weak-lensing plus substructure fit: the
  substructure likelihood uses the Meneghetti excess factor
  $f_{\rm GGSL}=3[1+(N-1)(1-\einf)]$, not Eq.~\ref{eq:integral}; the
  per-lens count appears in that script only in a printed cross-check.
  The best fit $N=10$, $k_t^{\rm app}=1.32$, $k_t^{\rm phys}=0.52\,h\,{\rm
  Mpc}^{-1}$, $\Delta\chi^2=-53.6$, the Limber analysis and the current
  $P(N)$ stand; \texttt{joint\_constraint.npz} and
  \texttt{lnL\_sub\_grid.npz} were not rewritten.
\item Paper VI's convergence profiles, the 38 percent convergence
  excess at the Einstein radius, the bound-mass table, the
  $\varepsilon$-versus-mass relation, the Meneghetti partial
  explanation (excess factor, not the integral), the \lcdm\ baselines
  of 0.43 and 0.05 per lens, the sub-galactic contribution below 1
  percent, and the scorecard's test outcomes other than the counts.
  Falsification criteria (2)--(6) are unchanged; criterion (1)'s
  significance falls from 10.4 to 8.8$\sigma$ and stays above 5.
\item Paper VIII's demographics: the 497 grade-A/B candidates, the
  7.88 per square degree rate, the 110,269 wide-survey lenses, the
  $\theta_E$ distribution and the efficiency models.
\item Paper IX entirely.
\item Paper X (Section~\ref{sec:paperx}).
\end{itemize}

\section{Falsification statements, restated}
\label{sec:falsification}

\textbf{Paper VI, criterion (1).} Absence of the $10^{10}$--$10^{12}\,\Msun$
Schechter excess in a sample of more than 1,000 lenses at 5 percent
efficiency would falsify ISF at more than $5\sigma$ (8.8$\sigma$ at
$N=6$; was 10.4).

\textbf{Paper VII.} If Euclid substructure counts at
$10^{8.8}\,\Msun$ yield fewer than 1,307 detections across 75,000
grade-A/B strong lenses (was 1,844), the ISF model is excluded at
$>95$ percent confidence for all $N\geq2$; fewer than 2,384 (was 3,358)
excludes $N\geq4$.

\textbf{Paper VIII.} If a complete perturber search of the 312
modelled Q1 grade-A/B systems at $10^{8.8}\,\Msun$ yields zero
detections, ISF with $N\geq6$ is excluded at $\geq4.2\sigma$ (was
$\geq5\sigma$) and all $N\geq2$ at $\geq2.9\sigma$ (was $\geq3.5$); if
Euclid DR1 substructure counts fall below 180--244 detections across
$\sim$15,000 grade-A/B lenses (was 257--348), all ISF models ($N\geq2$)
are excluded at $>95$ percent confidence; a DR1 null excludes even
$N=2$ at $\geq20\sigma$ (was $\geq24$).

\section{Consistency with Paper X}
\label{sec:paperx}

Paper X replaces Paper VI's intact-NFW perturber with the captured
spheroid its own chain derives and maps perturber mass to number with
a unit Jacobian. Its corrected per-lens count at $N=6$ with the extent
held, 0.756, sits beside the Jacobian-corrected 0.771 here: at $N=6$
Paper X's threshold change was this correction in disguise, the profile
effect being nil there. Its 312-system expectation for the registered
$m=1$ model with the extent held, 8.8--11.8, sits beside the corrected
9.0--12.0. Paper X's headline numbers are for the $m=2$--3 family and
are not affected; its concession~6 and its Table~3 row ``VIII, $N=6$''
(1.080; 16.8 / 13.7 / 12.6; 5.8 / 5.2 / 5.0) now read 0.771; 12.0 / 9.8 /
9.0; 4.9 / 4.4 / 4.2.

\section{Reproducibility}

The audit and the two reruns are three scripts of a few hundred lines,
each gated on reproducing the published product before it prints a
corrected number, published with the series' analysis code in the
\texttt{errata/} folder of the public repository together with the two
\texttt{\_v2} data products. The published Zenodo records of Papers VI
\citep{paper6}, VII \citep{paper7} and VIII \citep{paper8} receive new
versions carrying this note; the record numbers the series cites stay
valid.

\begin{thebibliography}{}
\bibitem[Baldry et al.(2012)]{baldry12} Baldry, I.~K., et al.\ 2012,
MNRAS, 421, 621
\bibitem[Merrill(2025)]{paper6} Merrill, C.~K.\ 2025, ISF Paper VI,
Galaxy-Mass Dark Substructure from Strong Lensing: A Parameter-Free Test of Independent Structure Formation, Zenodo,
doi:10.5281/zenodo.18637044
\bibitem[Merrill(2026a)]{paper7} Merrill, C.~K.\ 2026a, ISF Paper VII:
Joint Weak Lensing and Strong Lensing Substructure Constraints on ISF Cross-Branch Correlation Parameters, Zenodo,
doi:10.5281/zenodo.18644812
\bibitem[Merrill(2026b)]{paper8} Merrill, C.~K.\ 2026b, ISF Paper VIII:
Anchoring the Falsification Forecast to Measured Euclid Q1 Strong-Lens
Demographics, Zenodo, doi:10.5281/zenodo.21195841
\bibitem[Merrill(2026c)]{paper10} Merrill, C.~K.\ 2026c, ISF Paper X:
The Galaxy-Scale Dark Mass Profile Under Independent Structure
Formation, Zenodo, doi:10.5281/zenodo.23286778
\end{thebibliography}

\end{document}
